% Thesis summary in English, on the model of the summary of the reference thesis (3 A4 pages, no figures or tables).
% Every number comes from the thesis after commit 4cbefad (16 November 2026); see the chapters for the sources.
\documentclass[11pt,a4paper]{article}
\usepackage[utf8]{inputenc}
\usepackage[british]{babel}
\usepackage[a4paper,left=2.25cm,right=2.25cm,top=2.5cm,bottom=2.5cm,headheight=14pt]{geometry}
\usepackage{amsmath}
\usepackage{fancyhdr}
\usepackage[expansion=false]{microtype}  % no font expansion: body at a uniform 10.95 pt, as in the reference summary
\usepackage[hidelinks]{hyperref}

\hypersetup{pdftitle={Thesis Summary: Unsupervised Detection of Meteorological Anomalies with Spatio-Temporal Generative Adversarial Networks},
            pdfauthor={Alberto Pedalino}}

\pagestyle{fancy}
\fancyhf{}
\fancyhead[L]{Thesis Summary: Unsupervised Detection of Meteorological Anomalies}
\fancyhead[R]{Alberto Pedalino}
\fancyfoot[C]{\thepage}
\renewcommand{\headrulewidth}{0.4pt}

\newcommand{\Gn}{$G_{\mathrm{n}}$}

\begin{document}

\begin{center}
    {\Large\bfseries Thesis Summary\\[2pt]
    Unsupervised Detection of Meteorological Anomalies\\ with Spatio-Temporal Generative Adversarial Networks\par}
    \vspace{0.8em}
    {\large\textbf{Candidate:} Alberto Pedalino\par}
    \vspace{0.2em}
    {\large\textbf{Supervisor:} Simone Monaco\par}
    \vspace{0.2em}
    \textit{Politecnico di Torino}\\
    \textit{Master of Science in Computer Engineering}\\
    \textit{Software Engineering -- A.Y. 2025/2026}
\end{center}
\vspace{0.2em}

\section{Introduction and Motivation}

A meteorological anomaly is a state of the atmosphere that departs from what is usual for the place and the time of year. Reanalyses such as ERA5 describe the atmosphere on a regular grid over decades, but carry no anomaly labels: a detector must be unsupervised, and can only be evaluated indirectly. The work started from photovoltaic forecasting, where a forecaster is less reliable in rare conditions; the few variables of the PVGIS sites proved too limited to study them, hence the move to ERA5.

The thesis adapts STGAN, a spatio-temporal generative adversarial network proposed for road traffic, to gridded ERA5 fields. A generator predicts each field three hours ahead from the previous ones, and a discriminator learns to tell observed patches from reconstructed ones. The anomaly score adds two components: G, the generator's reconstruction error, and D, the difference between the discriminator's judgements of the observed and of the reconstructed patch.

The thesis asks four questions: which component to use; whether to add them, as the original score does; how to make the score comparable across places and seasons; and what the chosen score detects compared with simpler methods. The contribution is the adaptation of STGAN to unlabelled weather fields, the study of the two components of its score on these data, from which the choice of the detector score follows, and an evaluation without labels, against baselines and a reference built from the agreement of the methods; the building blocks already exist, and no new detection algorithm is proposed. In short, the detector uses the generator error because its meaning is clear; around the events on which most methods agree, it mostly flags rapid changes, as the three-hour change does, rather than states far from normal.

\section{Methodology}

\subsection{Detector and Score}
The generator is a convolutional recurrent network (ConvGRU) that reconstructs each field from the previous ones, three hours apart. Its inputs are standardised departures from the seasonal normal of each cell, day of year and hour, estimated on the training years. G is the reconstruction error, with each variable weighted inversely to how much it usually changes, averaged over 20 Monte Carlo dropout draws. The detector score, \Gn, expresses G relative to its own climatology for the same cell, day of year and hour, estimated on the validation years.

\subsection{Baselines}
\Gn{} is compared with eight baselines, each thresholded, like \Gn, at the 95th percentile of its own score. Seven read only one cell at one time step: six shallow models (PCA, CBLOF, MCD, HBOS, COPOD, Isolation Forest) and the distance from normal, the weighted squared departure of the cell. The eighth, the three-hour change, measures the same distance from the values three hours before. Flagged values in contact in space or in time form clusters.

\subsection{Concordant Reference}
The concordant reference is the set of values flagged by at least five of the nine methods, grouped into concordant clusters, short (lasting at most 72~h) or long. It is the agreement of the methods, not a ground truth. Each method is compared with it in four ways. The first is how many concordant clusters it detects: a detection needs one of its clusters within 300~km at a shared time step. The second and third are when its clusters start and end, and where they lie along the motion of the concordant cluster. The fourth is what its other flags nearby look like, also through quantities that no score reads, such as CAPE and the K index (convective energy and thunderstorm propensity).

\section{Experimental Setup}

The data are fifteen ERA5 variables, ten at the surface and five at 850~hPa, such as temperature, humidity, wind, pressure, cloud cover, solar radiation, precipitation and soil moisture. They are taken every three hours on 10\,611 cells of $0.5^\circ$ over Europe, the Mediterranean and parts of northern Africa and the Middle East. The years 1980--1999 are used for training, 2000--2004 for validation and 2005--2025 for testing. Among 29 trials of a hyperparameter search, the configuration was chosen by the maximum mean discrepancy (MMD) between observed and reconstructed validation fields. The model, its normalisations, the choice of the variables whose linear trend is removed and the thresholds use only the training and validation years; the 72~h split between short and long clusters is taken from the literature and not tuned on the test clusters.

\section{Results and Discussion}

\subsection{Choice of the Score}
On the validation years G and D order the cells almost without link: their Spearman correlation across the cells is on average 0.015. As in STGAN, the combined score adds the two components after normalising each, here with a min--max whose extremes are estimated on the validation years; it is in practice D, with a Spearman correlation of 0.99 and, with the three scores climatologically normalised, 94\% of its flags also flagged by D. G has a clear meaning, while D can be explained only in part; the detector therefore uses G, a choice that rests on its meaning, not on a better detection, which cannot be measured without labels.

\subsection{Comparability and Overlap with the Baselines}
With the climatological normalisation the share of flagged values stays between 3.7\% and 6.3\% in every latitude band and month on the validation years, against 1.6--11.1\% without it: the detector flags about as often in every place and season, as expected by construction. \Gn{} overlaps with the simple baselines only in part: only 41\% of its flags exceed the threshold of the distance from normal, and the closest baseline, the three-hour change, flags 54\% of them. This is partly expected, since \Gn{} measures how unexpected a state is three hours ahead, not how far it is from normal.

\subsection{Agreement with the Concordant Reference}
\Gn{} detects 46\% of the short concordant clusters and 91\% of the long ones; the criterion is loose, so that all methods detect 76--92\% of the long clusters. Of the clusters of \Gn, 54.5\% match no concordant cluster, but three in four of them lie next to concordant clusters too small to be studied. Its timing differs from that of the closest baselines. In short clusters MCD and the three-hour change start before \Gn{} in about half of the concordant clusters, and \Gn{} before them in one in five or fewer. In long clusters its flags end before those of PCA, MCD and the three-hour change and lie on average 19~km further along the mean motion of the concordant cluster, a fraction of a cell. This is partly expected for a reconstruction error, which is high where the field changes.

\subsection{Flags around the Concordant Clusters}
Around the concordant clusters, 90--96\% of the other flags of the single-cell methods have at least one variable beyond two standard deviations, against 40--48\% of the values that no method flags: these methods mark states far from normal. \Gn's flags more often follow rapid changes, like those of the three-hour change: the largest three-hour change before them is about twice that of the values no method flags. Around long clusters they also show unusual CAPE, K index and column water vapour more often than the flags of any other method (10.6--15.7\% against 5.2--6.6\%), by a small margin over MCD and, for CAPE, over the three-hour change. These results hold with stricter thresholds and other protocol values, except, in long clusters, the end against the other baselines and, at the 99th percentile, the lead on CAPE, K index and water vapour; with a minimum of 1125 values \Gn{} drops to fourth for detected long clusters.

\section{Conclusions}

Normalised and added as in STGAN, here with a min--max, the two components give in practice D, so the detector uses the generator error; its climatological normalisation makes the flag rate nearly uniform across places and seasons. Compared with simpler methods, \Gn{} overlaps with them only in part. Around the concordant clusters it mostly flags rapid changes, like the three-hour change and partly MCD, but it differs from both: its flags usually start later, in long clusters end earlier and lie further along the mean motion, and include moderate values that changed quickly but not enough to exceed the threshold of the three-hour change. Around long clusters, at the 95th and 97th percentiles, CAPE, K index and water vapour are unusual in its flags more often than for the other methods, by a small margin over MCD.

Beyond the choice of score, the work offers a way to study a detector when no labels exist: compare the components of a score before adding them, and describe what the detector flags against simple baselines and a concordant reference, even when its accuracy cannot be measured.

\end{document}
